Representing knowledge in structured forms that intelligent systems can query and reason
over is a foundational challenge of artificial intelligence. Knowledge graphs (KGs) have
become the attractive medium for this purpose, encoding entities, relationships, and semantic
descriptions extracted from text to support applications from question answering and
recommendation to predictive modelling~\cite{ji2021kgcompletion}. The ability to
construct KGs from unstructured text at scale has been substantially advanced by large
language models (LLMs), which enable flexible extraction of entities and relations without
hand-crafted rules~\cite{mo2025kggen}. Despite this progress, a fundamental tension
persists in KG construction: ontology and graph structure are typically specified
\emph{a priori} and evaluated independently of the tasks for which the graph is
ultimately used.

Existing KG construction tools --- from manual ontology engineering to LLM-driven
frameworks such as LlamaIndex and LangChain --- treat the ontology as a fixed design
choice established before construction begins and unchanged throughout.
Evaluation of the resulting graph is typically conducted through structural
quality metrics: completeness, consistency, and redundancy~\cite{paulheim2017kgquality}.
\citet{paulheim2024structuralquality} draw a sharp distinction between this
intrinsic mode of evaluation and task-based evaluation --- assessing a
KG by its utility on a concrete downstream application. Frameworks such as
KGrEaT~\cite{heist2023kgreat} and KGBench~\cite{bloem2021kgbench} try to formalise the approach to such evaluations.
As \citet{heist2023kgreat} observe: \textit{``Performance indicators may
be used to make an informed decision when picking a KG for a given task\ldots\ and to
compare different versions of a KG during its life cycle.''} Yet despite this insight,
no existing work closes the loop at construction time: no downstream evaluation signal has been used to directly steer KG ontology evolution during the
building process itself

This paper aims to address the gap between graph construction and downstream task
performance by introducing a feedback-driven framework in which the KG
ontology is treated as a tunable hyperparameter. Rather than fixing the schema in
advance, we propose a closed-loop system:
\begin{enumerate}
    \item[\textit{(i)}] An Ontology Agent proposes candidate ontology configurations.
    \item[\textit{(ii)}] A KG Builder Agent constructs knowledge graphs from text under each configuration.
    \item[\textit{(iii)}] Graph-derived features are computed and evaluated on a downstream task.
    \item[\textit{(iv)}] The resulting task performance score is returned as a feedback signal to guide the next ontology revision.
\end{enumerate}
This loop repeats until performance converges.

We instantiate this framework on the task of next-day stock price movement prediction
using the FNSPID financial news dataset~\cite{dong2024fnspid}, a corpus of financial
news articles linked to NASDAQ price time series. This domain offers a clear and
measurable downstream objective, dense entity interactions, and a natural temporal
structure that motivates subgraph-based feature extraction. The framework is
designed to be domain-agnostic and is extensible to other settings.

\paragraph{Contributions.} This paper makes the following contributions:
\begin{enumerate}
    \item We propose a feedback-driven KG construction framework that treats ontology
    as a hyperparameter optimized for downstream task performance, closing the loop
    between graph construction and task evaluation.

    \item We introduce a structured feature representation for temporal KG subgraphs
    --- comprising node type histograms, relation type histograms, metapath counts,
    temporal novelty statistics, and topology features --- that supports lightweight
    downstream task evaluation without requiring graph neural networks.

    \item We provide empirical evidence on the FNSPID dataset that task-aware ontology
    adaptation improves downstream task performance, enables more compact graph representations
    without sacrificing performance, and produces more informative graph structures
    over successive iterations.
\end{enumerate}

\paragraph{Paper organisation.} Section~\ref{sec:related} reviews related work on KG
construction, task-based KG evaluation, and task-aware graph learning.
Section~\ref{sec:hypotheses} states the research hypotheses formally.
Section~\ref{sec:methodology} describes the proposed framework in detail.
Section~\ref{sec:experiments} presents the experimental setup, and
Section~\ref{sec:results} reports results. Section~\ref{sec:discussion} discusses
findings and limitations, and Section~\ref{sec:conclusion} concludes.
